% Gabarit pandoc des corrigés publiés (brut et guidé). Contenu sans marque ni nom de site.
% Compilation : lualatex (latexmk -lualatex).
% Le Markdown reste la source unique : ne pas éditer ce fichier .tex à la main.
\documentclass[11pt,a4paper]{article}

% --- Moteur, polices (Latin Modern) et langue ---------------------------------
\usepackage{iftex}
\ifLuaTeX\else\errmessage{Ce document se compile avec lualatex}\fi
\usepackage{amsmath,amssymb,mathtools}
\usepackage{fontspec}
% mathrm=sym : \mathrm prend ses glyphes dans la police mathématique (ℓ, μ, Ω), comme KaTeX.
\usepackage[math-style=TeX,bold-style=TeX,mathrm=sym]{unicode-math}
% Repli pour les glyphes absents de Latin Modern (indices et exposants Unicode,
% flèches, filets des schémas en texte, symboles) : DejaVu, puis Latin Modern Math
% (ℓ, lettres mathématiques Unicode des légendes).
\directlua{
  luaotfload.add_fallback("repliTexte", {"DejaVuSerif:mode=harf;", "DejaVuSans:mode=harf;", "file:latinmodern-math.otf:mode=harf;"})
  luaotfload.add_fallback("repliMono", {"DejaVuSansMono:mode=harf;", "DejaVuSans:mode=harf;", "file:latinmodern-math.otf:mode=harf;"})
}
\setmainfont{Latin Modern Roman}[RawFeature={fallback=repliTexte}]
\setsansfont{Latin Modern Sans}[RawFeature={fallback=repliTexte}]
\setmonofont{Latin Modern Mono}[RawFeature={fallback=repliMono}]
\setmathfont{Latin Modern Math}
\usepackage[french]{babel}
\frenchsetup{SmallCapsFigTabCaptions=false}

% --- Mise en page --------------------------------------------------------------
\usepackage[a4paper,margin=2.2cm,headheight=14pt,footskip=1.2cm]{geometry}
\usepackage{microtype}
\setlength{\emergencystretch}{4em}
\tolerance=1500
\setlength{\parindent}{0pt}
\setlength{\parskip}{0.55\baselineskip plus 2pt minus 1pt}
\usepackage{xcolor}
% Gris nommés : babel-french rend « ! » actif dans le document.
\definecolor{gris2}{gray}{0.98}
\definecolor{gris3}{gray}{0.97}
\definecolor{gris5}{gray}{0.95}
\definecolor{gris35}{gray}{0.65}
\definecolor{gris40}{gray}{0.60}
\definecolor{gris45}{gray}{0.55}
\definecolor{gris50}{gray}{0.50}
\definecolor{gris55}{gray}{0.45}
\definecolor{gris60}{gray}{0.40}
\definecolor{gris65}{gray}{0.35}
\definecolor{gris70}{gray}{0.30}
\definecolor{gris80}{gray}{0.20}
\definecolor{gris85}{gray}{0.15}
\usepackage{graphicx}
\usepackage{float}
\usepackage{caption}
\captionsetup{font=small,labelformat=empty,justification=centering}
\usepackage{booktabs,longtable,array,multirow,calc}
\usepackage{siunitx}
\sisetup{locale=FR}
\usepackage[version=4]{mhchem}
\usepackage{listings}
\usepackage{fancyhdr}
\usepackage{lastpage}
\usepackage[most]{tcolorbox}
\usepackage{enumitem}
\setlist{nosep,leftmargin=*,topsep=2pt}
\usepackage[hidelinks,unicode]{hyperref}
\usepackage{bookmark}
\hypersetup{pdftitle={Agro-Véto TB — Anglais (LVO) 2026},pdfauthor={Rédigé par IA},pdfcreator={LaTeX},pdfsubject={Corrigé brut}}
\usepackage{xurl}

% --- Pandoc : listes, tableaux, figures ----------------------------------------
\providecommand{\passthrough}[1]{#1}
\providecommand{\tightlist}{\setlength{\itemsep}{0pt}\setlength{\parskip}{0pt}}

\makeatletter
\def\maxwidth{\ifdim\Gin@nat@width>\linewidth\linewidth\else\Gin@nat@width\fi}
\def\maxheight{\ifdim\Gin@nat@height>0.55\textheight 0.55\textheight\else\Gin@nat@height\fi}
\makeatother
\setkeys{Gin}{width=\maxwidth,height=\maxheight,keepaspectratio}
\setcounter{secnumdepth}{-\maxdimen}
\providecommand{\pandocbounded}[1]{#1}

% --- Code (listings) -----------------------------------------------------------
\lstset{
  basicstyle=\ttfamily\small,
  columns=fullflexible,
  keepspaces=true,
  breaklines=true,
  breakatwhitespace=false,
  postbreak=\mbox{\textcolor{gris50}{$\hookrightarrow$}\space},
  frame=single,
  rulecolor=\color{gris35},
  framesep=4pt,
  xleftmargin=5pt,xrightmargin=5pt,
  backgroundcolor=\color{gris3},
  keywordstyle=\bfseries,
  commentstyle=\color{gris65},
  stringstyle=\color{gris80},
  showstringspaces=false,
  tabsize=4,
  extendedchars=true,
  aboveskip=6pt,belowskip=6pt,
}

% --- Macros des corrigés (KaTeX ↔ LaTeX) ---------------------------------------
% Les corrigés sont écrits pour KaTeX : on définit ici ce que KaTeX connaît et
% que LaTeX n'a pas, sans jamais retoucher le Markdown.
\providecommand{\R}{\mathbb{R}}
\providecommand{\N}{\mathbb{N}}
\providecommand{\Z}{\mathbb{Z}}
\providecommand{\Q}{\mathbb{Q}}
\providecommand{\C}{\mathbb{C}}
\providecommand{\dd}{\mathrm{d}}
\providecommand{\llbracket}{\lBrack}
\providecommand{\rrbracket}{\rBrack}
\providecommand{\lBrace}{\lbrace}
\providecommand{\rBrace}{\rbrace}
\providecommand{\Box}{\square}
\providecommand{\Diamond}{\lozenge}
\providecommand{\varnothing}{\emptyset}
\providecommand{\argmax}{\operatorname*{arg\,max}}
\providecommand{\argmin}{\operatorname*{arg\,min}}
\renewcommand{\leq}{\leqslant}
\renewcommand{\geq}{\geqslant}

% Équation hors texte : réduite à la largeur de la ligne si elle déborde.
\newsavebox{\boiteequation}
\newcommand{\equationajustee}[1]{%
  \sbox{\boiteequation}{$\displaystyle #1$}%
  \ifdim\wd\boiteequation>\linewidth
    \[\resizebox{\linewidth}{!}{\usebox{\boiteequation}}\]%
  \else
    \[#1\]%
  \fi}

% --- Encadrés pédagogiques (sobres, lisibles en noir et blanc) -----------------
\tcbset{
  encadre/.style={
    enhanced,breakable,sharp corners,
    colback=white,colframe=gris55,
    boxrule=0.4pt,leftrule=2.2pt,
    left=6pt,right=6pt,top=3pt,bottom=3pt,
    before skip=6pt,after skip=6pt,
    fonttitle=\bfseries\small,coltitle=black,
    attach title to upper,after title={\quad},
    parbox=false,
  }
}
\newtcolorbox{indication}[1]{encadre,colframe=gris45,colback=gris2,title={Indication~#1}}
\newtcolorbox{methode}{encadre,colframe=gris70,title={Méthode}}
\newtcolorbox{piege}{encadre,colframe=gris85,boxrule=0.8pt,leftrule=2.2pt,title={Piège}}
\newtcolorbox{jury}{encadre,colframe=gris60,colback=gris5,title={Jury}}
\newtcolorbox{encadrenote}{encadre,colframe=gris40}

% --- Titres de section ---------------------------------------------------------
\usepackage{titlesec}
\titleformat{\section}{\Large\bfseries}{}{0pt}{}
\titleformat{\subsection}{\large\bfseries}{}{0pt}{}
\titleformat{\subsubsection}{\normalsize\bfseries}{}{0pt}{}
\titlespacing*{\section}{0pt}{14pt plus 4pt minus 2pt}{6pt}
\titlespacing*{\subsection}{0pt}{10pt plus 3pt minus 2pt}{4pt}
\titlespacing*{\subsubsection}{0pt}{8pt plus 2pt minus 2pt}{2pt}

% --- En-tête et pied de page ---------------------------------------------------
\pagestyle{fancy}
\fancyhf{}
\fancyhead[L]{\small\textsc{Corrigé brut} · Agro-Véto TB}
\fancyhead[R]{\small Anglais 2026}
\fancyfoot[L]{\footnotesize Corrigé rédigé par IA, non officiel}
\fancyfoot[R]{\footnotesize \thepage\,/\,\pageref*{LastPage}}
\renewcommand{\headrulewidth}{0.4pt}
\renewcommand{\footrulewidth}{0pt}
\fancypagestyle{plain}{\fancyhf{}\fancyfoot[R]{\footnotesize \thepage\,/\,\pageref*{LastPage}}\fancyfoot[L]{\footnotesize Corrigé rédigé par IA, non officiel}\renewcommand{\headrulewidth}{0pt}}

\begin{document}
\thispagestyle{plain}
\begin{center}
  {\small\scshape Corrigé brut}\par\smallskip
  {\LARGE\bfseries Agro-Véto TB — Anglais (LVO) 2026\par}\medskip
  {\large Agro-Véto · TB · Anglais · 2026\par}
\end{center}
\vspace{2pt}\hrule\vspace{8pt}

\emph{Corrigé rédigé par IA, non officiel. Signalez toute erreur.}

\textbf{Problématique commune} : les deux articles (\emph{The New York Times}, 20 janvier 2026 ; \emph{Ici}, France Bleu, 10 février 2026) montrent le vin pris dans les tensions entre les États-Unis et la France, l'un sous l'angle de la menace diplomatique, l'autre sous celui de la riposte humoristique d'un vigneron. La partie I compare ce que disent les deux textes des effets de ces tensions ; la partie II demande une argumentation personnelle sur l'un des deux sujets au choix, que l'on traite ici tous les deux.

Durée : 2 heures, sans document ni appareil. Barème : partie I sur 10, partie II sur 10. Longueur imposée pour chaque partie : 220 mots ± 10 \%, soit 198 à 242 mots, nombre de mots indiqué à la fin. Le candidat ne traite qu'un des deux sujets d'expression et en indique le numéro. Nombre de mots : suites de caractères séparées par des espaces (mot à trait d'union = un mot, nombre = un mot, ponctuation seule non comptée).

\hypertarget{i-compruxe9hension-de-luxe9crit}{%
\section{I --- Compréhension de l'écrit}\label{i-compruxe9hension-de-luxe9crit}}

\hypertarget{how-do-political-tensions-affect-the-wine-sector-according-to-the-two-articles-answer-the-question-in-your-own-words.}{%
\subsection{How do political tensions affect the wine sector according to the two articles? Answer the question in your own words.}\label{how-do-political-tensions-affect-the-wine-sector-according-to-the-two-articles-answer-the-question-in-your-own-words.}}

\emph{Problématique : le vin subit les tensions politiques comme moyen de pression et source d'incertitude, mais peut aussi s'en servir comme outil de communication ; les deux articles l'abordent l'un par la menace, l'autre par l'humour.}

As both articles show, political tensions affect the wine sector in several ways, although the two texts adopt contrasting angles: The New York Times (20 January 2026) reports a diplomatic threat, while Ici (France Bleu, 10 February 2026) tells a light-hearted story.

In The New York Times, wine appears as a weapon. Donald Trump threatened 200 per cent tariffs on French wine and Champagne if France refused to join his ``Board of Peace''. A similar threat in March 2025 was never carried out, but it worried Champagne producers, whose largest foreign market is the United States. Producers thus live in uncertainty over unrelated disputes.

The France Bleu report shows that the consequences are lasting. The 200 or so winegrowers of the Costières de Nîmes appellation export part of their wine to the United States, where it is now taxed at 10 to 15 per cent. Although their sales there rose slightly last year, professionals fear that higher prices could turn American customers away from European wines.

Yet tensions can also become an opportunity. A winegrower from the Gard named his rosé ``For sure'', after Emmanuel Macron's viral expression at Davos, and offered it to the President, hoping to revive rosé sales in the United States: orders came in immediately.

Both articles thus present wine as a symbolic product caught up in politics: a means of pressure for Washington, but a means of communication for French producers.

\textbf{(236 words)}

\hypertarget{ii-expression-uxe9crite}{%
\section{II --- Expression écrite}\label{ii-expression-uxe9crite}}

\hypertarget{sujet-1-why-do-certain-agricultural-products-such-as-wine-wheat-or-dairy-products-become-political-or-symbolic-tools-in-international-relations-give-examples.}{%
\subsection{Sujet 1 --- Why do certain agricultural products (such as wine, wheat, or dairy products) become political or symbolic tools in international relations? Give examples.}\label{sujet-1-why-do-certain-agricultural-products-such-as-wine-wheat-or-dairy-products-become-political-or-symbolic-tools-in-international-relations-give-examples.}}

\emph{Problématique : certains produits agricoles deviennent des armes parce qu'ils sont à la fois vitaux, chargés d'identité nationale et portés par des producteurs qui pèsent sur la vie politique.}

\textbf{Subject 1}

In January 2026, Donald Trump threatened to tax French wine at 200 per cent because France had refused to join his ``Board of Peace''. Why do certain agricultural products become political or symbolic tools in international relations?

Not all products carry the same weight. First, some are vital, which makes them powerful levers. Wheat feeds whole populations, so controlling its supply gives a country leverage. After invading Ukraine in 2022, Russia blocked Ukrainian grain exports through the Black Sea, and food prices rose in many importing countries. Food thus becomes a weapon.

Secondly, other products are symbols of national identity. Champagne, Roquefort or Scotch whisky embody a country's know-how and are protected by designations of origin. Hitting them hurts a nation's pride as much as its economy. This is why the United States imposed a 100 per cent duty on Roquefort in 1999 during the dispute over hormone-treated beef, and taxed French wine in 2019 during the Airbus-Boeing conflict.

Thirdly, agricultural products affect voters. Farmers are well organised in many countries, so targeting their exports puts indirect pressure on governments. When China taxed Australian wine heavily in 2020, after Canberra had called for an inquiry into the origins of Covid-19, Australian winemakers lost their biggest export market.

Agricultural products therefore become tools because they combine economic importance, emotional value and political weight. Yet such pressure seldom changes a government's policy and mainly hurts producers.

\textbf{(234 words)}

\hypertarget{sujet-2-do-you-think-social-media-and-communication-strategies-can-change-the-impact-of-political-tensions-give-examples.}{%
\subsection{Sujet 2 --- Do you think social media and communication strategies can change the impact of political tensions? Give examples.}\label{sujet-2-do-you-think-social-media-and-communication-strategies-can-change-the-impact-of-political-tensions-give-examples.}}

\emph{Problématique : les réseaux sociaux amplifient les tensions et la communication peut en atténuer les effets, mais ni l'un ni l'autre ne change les décisions politiques.}

\textbf{Subject 2}

When Donald Trump threatened French wine with 200 per cent tariffs, a winegrower from the Gard replied with a rosé called ``For sure'', inspired by Emmanuel Macron's viral words at Davos. Can social media and communication strategies change the impact of political tensions?

To begin with, social media amplify tensions. Political leaders now announce sanctions or tariffs directly online, as Donald Trump often does on Truth Social. Threats therefore spread instantly and worry markets and producers before any decision has been taken. A diplomatic quarrel can quickly turn into a boycott.

Nevertheless, communication can also soften the effects of a crisis, or even turn it into an opportunity. The ``For sure'' label made a tariff dispute amusing and immediately brought in orders. Similarly, in 2022, Volodymyr Zelensky's daily videos helped Ukraine win international sympathy and military aid. In Canada, the ``Buy Canadian'' movement that spread online in 2025 in response to American tariffs encouraged consumers to support local producers.

However, the influence of these strategies remains limited. A viral bottle does not remove customs duties: Costières wines still pay 10 to 15 per cent to enter the United States. Communication changes the way tensions are perceived, not the political decisions themselves.

In conclusion, I believe that social media and communication strategies can change the impact of political tensions, mainly by shaping perceptions and creating new markets, but they cannot replace diplomacy.

\textbf{(231 words)}

\end{document}
